\section{Introduction}
\subsection{Background and Rationale}
Cooperative survey coverage divides an area among vehicles so they can observe it in parallel. Coverage path planning includes marine applications, but coverage of an ideal grid does not by itself establish bathymetric reconstruction quality \cite{galceran2013coverage,yang2023cooperative}. Marine coordination also connects communication quality with task allocation \cite{martorell2024coordination}. A fixed partition is a useful reference when workloads and progress are balanced. Online reassignment may become useful after an unexpected slowdown, provided the vehicles can exchange current coverage and task state.

Our motivating problem is whether reserving more work before a predicted outage improves a two-vehicle survey. Prediction may help a vehicle avoid exhausting its committed work, but an always-long queue could provide the same benefit. Moreover, route changes can alter connectivity, and abrupt synthetic fades may have no observable precursor. The experiment therefore retains both a communication-independent reference and a non-predictive long-queue control.

\subsection{Theory and Related Work}
Yang et al. address cooperative USV coverage planning \cite{yang2023cooperative}; Bramblett and Bezzo study multi-robot planning with restricted communication \cite{bramblett2023epistemic}. These motivate a bounded coordination experiment rather than a new vessel model. Published consensus-based auctions provide an allocation reference \cite{choi2009auction}, but our deterministic greedy rule does not implement CBBA's distributed consensus or inherit its guarantees.

The model combines a log-distance signal form \cite{goldsmith2005pathloss} with a two-state burst process motivated by Gilbert \cite{gilbert1960burst}. Persistence and ordinary least squares are standard forecasting components \cite{hyndman2021forecasting,nistLeastSquares}. The combined residual-signal and relative-motion warning is a study-specific construction. We compare reservation policies under one common allocator, without claiming superiority to published coverage or allocation systems.

\subsection{Research Problem, Aim, and Research Questions}
The aim is to evaluate whether short-horizon outage warnings improve online task reservation between two USVs. The main question is: how does prediction-triggered reservation affect mission completion time and redundant coverage relative to reactive resynchronization under intermittent boat-to-boat connectivity?
\begin{enumerate}[label=\textbf{RQ\arabic*.}]
\item Across fade-persistence and balanced/unequal-progress conditions, what is the effect of trend-triggered reservation relative to reactive coordination?
\item Does any benefit remain relative to an always-long queue, and what communication and computation costs accompany it?
\item How useful are trend and persistence warnings in terms of false alarms, missed outages, and lead time under the same reservation action?
\end{enumerate}
The expectation that anticipation might help requires both a useful warning and an action that improves the mission. Null and adverse effects are valid outcomes. Fixed partitioning is a reference, not an assumed defective baseline.
